\documentclass[11pt]{article}
\usepackage[margin=0.8in]{geometry}
\usepackage{enumitem}
\usepackage{setspace}
\setlist[itemize]{leftmargin=1.15em, itemsep=0.1em, topsep=0.18em}
\setlist[enumerate]{leftmargin=1.35em, itemsep=0.12em, topsep=0.2em}
\setlength{\parindent}{0pt}
\setlength{\parskip}{0.28em}
\setstretch{1.05}
\pagestyle{plain}

\begin{document}

\begin{center}
{\Large\bfseries Principles of Microeconomics}\\[0.3em]
Sample syllabus\\[0.15em]
Ke Cheng HSU\\
Department of Economics, National Chengchi University\\
Undergraduate, 3 credits, one semester, taught in English
\end{center}

\textbf{Place in the curriculum.}
This is a one-semester course in graphs and algebra for students who have not taken intermediate theory. Calculus is not used. It can serve as the microeconomics half of a first-year introduction. The department's required freshman course, Economics, runs for a full year and also covers macroeconomics. Students who continue to sophomore Microeconomics meet consumer theory and producer theory with calculus in that course.

\textbf{Course description.}
The course studies how consumers and firms make choices, how prices are formed, and when a market outcome is efficient. Each week has one main graph or calculation. Lectures follow N.\ Gregory Mankiw, \emph{Principles of Microeconomics}, 10th edition (Cengage, 2024), with two chapters omitted. In that edition the costs of production are Chapter 14. Chapter 12, The Economics of Healthcare, and Chapter 13, The Design of the Tax System, are omitted because they are tied to U.S.\ institutions. Problem sets use the Questions for Review and the Problems and Applications at the end of the assigned chapter.

\textbf{How the course runs.}
There are three hours of class a week. The first two hours introduce the chapter. The third hour is a problem session: students work a new exercise in pairs, then one pair puts the graph on the board. Six problem sets are due at the start of class, about every two teaching weeks. Students may discuss the method. Each student writes and submits their own answers. Solutions are posted after the deadline. Two short quizzes reuse the form of a problem-set question, so the sets are the way to prepare. There is no make-up quiz. A documented absence drops that quiz. Office hours are for a graph or a calculation that did not come out in the problem session.

\textbf{Learning outcomes.}
By the end of the semester, students should be able to
\begin{itemize}
\item explain scarcity, opportunity cost, and the gains from trade with a production-possibilities frontier,
\item shift supply or demand and find the new equilibrium price and quantity,
\item compute price elasticity and say whether a price rise raises or lowers expenditure,
\item mark consumer surplus, producer surplus, and deadweight loss under a tax, a price ceiling, a tariff, and a monopoly,
\item draw average and marginal cost, find a competitive firm's shutdown price, and set marginal revenue equal to marginal cost for a monopolist.
\end{itemize}

\textbf{Assessment.}
Problem sets 20 percent. The lowest set is dropped. Quizzes 10 percent. Midterm examination 30 percent, closed book, covering weeks 1--8. Final examination 40 percent, closed book and cumulative for this course. A calculator with no stored notes is allowed on exams.

\textbf{Weekly plan.}
\begin{enumerate}
\item Ten principles, and models as simplified maps. Chapters 1--2. The production-possibilities frontier and a positive statement versus a normative one.
\item Gains from trade. Chapter 3. Comparative advantage from opportunity cost, and why both sides can gain when each specializes.
\item Supply, demand, and equilibrium. Chapter 4. Movements along a curve versus shifts. A shortage and a surplus.
\item Elasticity. Chapter 5. The midpoint formula. Elastic, inelastic, and unit-elastic demand. Income and cross-price elasticity in one application.
\item Price controls and taxes. Chapter 6. A binding price ceiling and a binding price floor. Tax incidence when demand is more inelastic than supply.
\item Consumer surplus, producer surplus, and efficiency. Chapter 7. The surplus triangles. A price change that is a transfer and a price change that destroys surplus.
\item The costs of taxation. Chapter 8. The deadweight-loss triangle. Why the loss rises faster than the tax.
\item International trade. Chapter 9. The world price, imports, exports, a tariff, and who gains inside the country.
\item Midterm examination on weeks 1--8.
\item Externalities. Chapter 10. A negative externality and a Pigouvian tax. The Coase theorem when bargaining is cheap, and one case where it is not.
\item Public goods and common resources. Chapter 11. Nonrival and nonexcludable goods. Free riding. A common resource that is overused.
\item The costs of production. Chapter 14. Opportunity cost versus accounting cost. Fixed and variable cost. Why marginal cost eventually rises. Economies of scale.
\item Firms in competitive markets. Chapter 15. Price equals marginal cost. The shutdown rule in the short run. Entry and exit in the long run, and a horizontal long-run supply in a constant-cost industry.
\item Monopoly. Chapter 16. Marginal revenue below the demand curve. The quantity where marginal revenue equals marginal cost, and the resulting deadweight loss. A natural monopoly and one regulatory rule.
\item Monopolistic competition. Chapter 17. Downward-sloping demand and free entry. Excess capacity in the long-run equilibrium.
\item Oligopoly. Chapter 18. A prisoners' dilemma for two price-setting firms. Why a cartel is profitable and why a member wants to cheat.
\item Factor markets, and a review. Chapter 19. A competitive firm's demand for labor as the value of the marginal product. The review returns to four graphs: equilibrium, a tax, costs, and monopoly.
\item Final examination.
\end{enumerate}

\end{document}
